\section{Conclusion and Future Work}
\label{sec:conclusion}

Across two model families and four reasoner scales, what governs document-agent accuracy on long, visually dense documents is not the size of the reasoner but whether it can actively control perception: an agent that writes code to crop, zoom, and query a frozen vision--language model on demand matches or exceeds much larger frozen models and far surpasses fixed-granularity tool agents. The advantage is mechanistically located --- it concentrates where cropping recovers detail a whole-page read misses, collapses when visual perception is replaced by optical character recognition, and appears only once the base model is capable enough, in both reasoning and code, to drive the loop.

The two strongest scaffolds in this family tie at the 27B backbone, and the append-only member preserves the growing-prefix trajectory that established weight-level training assumes; it is therefore the principled target for training a small agent. Our preliminary supervised study at the $\leq$8B tier is an honest first step --- a rejection-sampling fine-tuning pipeline and data pool whose quantitative verdict awaits a corrected evaluation. Training a small agent to realize the active-perception advantage, with on-policy methods that match the deployment distribution, is the open direction we point to.
